\section{Conclusiones}
Los ensayos de seguridad eléctrica efectuados en el marco de la norma IEC 60601-1 demostraron que el equipo evaluado posee un aislamiento adecuado entre la red de alimentación y las partes aplicadas al paciente, registrando valores de resistencia de aislamiento superiores a $100\text{ M}\Omega$ entre la tierra de protección y las partes aplicadas, así como líneas abiertas entre la red de alimentación y la tierra o las partes aplicadas. Asimismo, las corrientes de fuga se mantuvieron estrictamente por debajo de los límites normativos tanto en condición normal como en condición de fallo, registrándose corrientes de fuga a tierra inferiores a $30\ \mu\text{A}$ en operación normal y corrientes de fuga en partes aplicadas de hasta $0{,}8\ \mu\text{A}$. Sin embargo, la medición de resistencia a tierra de protección arrojó un valor neto de $0{,}500\ \Omega$, el cual supera significativamente el límite máximo permitido por la normativa ($0{,}100\text{--}0{,}200\ \Omega$), evidenciando un fallo en la continuidad del chasis que requiere revisión técnica obligatoria antes de autorizar su uso en servicio.

En lo relativo a la evaluación de eficacia funcional, tanto el monitor multiparamétrico como el oxímetro de pulso exhibieron un excelente desempeño y una elevada exactitud frente a las señales de referencia generadas por los simuladores Fluke ProSim 4 y Fluke ProSim SPOT Light. Las desviaciones registradas en la frecuencia cardíaca resultaron ser de a lo sumo $1\text{ bpm}$ (error de $0{,}83\%$ a $120\text{ bpm}$), mientras que los errores en la medición de la saturación de oxígeno ($\text{SpO}_2$) no superaron el $1\%$, manteniéndose holgadamente dentro de las tolerancias clínicas aceptables del $\pm 2\%$. De igual forma, el sistema de alarmas fisiológicas respondió de manera precisa y oportuna ante la simulación de condiciones críticas de bradicardia y taquicardia.

En función de los resultados obtenidos, se concluye que el equipamiento satisface los criterios de eficacia y exactitud operacional para la monitorización de signos vitales, pero no resulta completamente apto para su aplicación en un entorno clínico real hasta tanto no se lleve a cabo el mantenimiento correctivo correspondiente para subsanar la elevada resistencia de la línea de tierra.